% !TEX program = lualatex
\documentclass[aspectratio=43,professionalfonts,handout]{beamer}
\usepackage{beamer_template}

\newcommand{\LectureNum}{15}
\newcommand{\LectureTitle}{前期復習}
\newcommand{\TextPages}{p.\,全範囲（pp.46-103）}
\newcommand{\CourseName}{応用数学}
\renewcommand{\TermName}{前期}

\begin{document}

\begin{frame}{\CourseName\quad 第\LectureNum 回「\LectureTitle」}
  {\small \TermName\quad \TeacherName\quad ／\quad 教科書 \TextPages}

  \textbf{目標}：第1〜14週の内容を総復習し、前期末試験に向けて理解を確認する

\end{frame}

\begin{frame}{ラプラス変換まとめ}
  \begin{block}{}
  \begin{itemize}
    \item 定義・基本公式一覧
    \item 各法則の使い分け
    \item 微分方程式の解法手順
  \end{itemize}
  \end{block}
\end{frame}

\begin{frame}{フーリエ級数まとめ}
  \begin{block}{}
  \begin{itemize}
    \item フーリエ係数の公式（周期 $2\pi$ ・ 2l ）
    \item 偶関数・奇関数の場合
    \item 複素フーリエ係数
    \item 収束定理
  \end{itemize}
  \end{block}
\end{frame}

\begin{frame}{フーリエ変換まとめ}
  \begin{block}{}
  \begin{itemize}
    \item フーリエ変換の定義と反転公式
    \item 6つの性質
    \item たたみ込みとの関係
    \item スペクトルの概念
  \end{itemize}
  \end{block}
\end{frame}

\begin{frame}{総合演習}
  \begin{exampleblock}{自分で解いてみよう}
  \begin{itemize}
    \item 各章の練習問題から重要問題を選題
    \item 試験頻出パターンの確認
  \end{itemize}
  \end{exampleblock}
\end{frame}

\begin{frame}{前期末に向けて}
  \begin{itemize}
    \item 評価：中間35%・期末35%・課題30%
    \item 60点以上で合格
    \item ラプラス変換表・フーリエ変換表の活用
  \end{itemize}
\end{frame}

\end{document}
